\chapter{Ordliste}

\begin{description}[style=nextline,leftmargin=0pt,labelindent=0pt]
  \item[Adækvans] En begrænsning af ansvar til følger, der er tilstrækkeligt nært forbundet med handlingen og ikke helt ekstraordinære.
  \item[Anbringende] En juridisk eller faktisk grund, som en part bruger til at støtte sin påstand.
  \item[Appel] Prøvelse af en afgørelse ved en højere instans efter regler om frist og adgang.
  \item[Afgørelse] En beslutning, som en offentlig myndighed træffer om en sag med retlig virkning for en eller flere personer.
  \item[Begrundelse] Den del af en afgørelse eller dom, der forklarer retsgrundlag, faktum og afgørende hensyn.
  \item[Bevisbyrde] Risikoen for, at et omstridt faktum ikke lægges til grund.
  \item[Besiddelse] Faktisk kontrol over en ting eller et gode.
  \item[Bindende regel] En regel, som adressaten skal følge, og som kan håndhæves gennem en retsfølge.
  \item[Charteret] Den Europæiske Unions charter om grundlæggende rettigheder.
  \item[Culpa] Et ansvarsgrundlag, der bygger på uforsvarlig adfærd, typisk uagtsomhed; forsæt kan også opfylde skyldkravet.
  \item[Databehandler] En aktør, der behandler personoplysninger på vegne af den dataansvarlige.
  \item[Dataansvarlig] Den, der bestemmer formålene med og hjælpemidlerne til behandling af personoplysninger.
  \item[Direktiv] En EU-retsakt, der binder medlemsstaterne med hensyn til det resultat, der skal nås.
  \item[Dom] En domstols afgørelse, normalt med resultat og begrundelse.
  \item[Ejendomsret] En samling af retlige beføjelser over et gode, eksempelvis brug, udelukkelse og salg.
  \item[EU-forordning] En EU-retsakt, der som udgangspunkt gælder direkte i medlemsstaterne.
  \item[Forarbejder] Dokumenter fra lovgivningsprocessen, som kan belyse formål og forståelse.
  \item[Forfatning] De grundlæggende regler om statens organisation, magt og rettigheder.
  \item[Fuldbyrdelse] Gennemførelse af en dom eller et andet tvangsgrundlag.
  \item[Fuldmagt] Beføjelse til at handle på en andens vegne.
  \item[Fortolkning] Fastlæggelse af, hvad en retlig tekst eller regel betyder.
  \item[Hjemmel] Den regel, der giver en myndighed eller person retlig adgang til at handle.
  \item[Inhabilitet] Forhold, der gør, at en person ikke må medvirke på grund af tvivl om upartiskhed.
  \item[Indgreb] En handling, der begrænser en rettighed eller påfører en person en byrde.
  \item[Kontradiktion] Muligheden for at kende og kommentere materiale, der kan få betydning for afgørelsen.
  \item[Legalitetsprincip] Kravet om retligt grundlag og forudsigelighed for offentlige indgreb; i strafferetten især ingen straf uden lov.
  \item[Legitimation] Det ydre grundlag for, at en tredjemand kan stole på, at en person har en beføjelse.
  \item[Mangelsbeføjelse] Et krav ved en ydelse, der ikke svarer til aftalen, eksempelvis afhjælpning, afslag eller ophævelse.
  \item[Mediation] En frivillig konfliktløsningsproces, hvor en neutral tredjeperson hjælper parterne med at finde en løsning.
  \item[Misligholdelse] Manglende eller mangelfuld opfyldelse af en aftale.
  \item[Officialprincippet] Myndighedens ansvar for, at en sag er tilstrækkeligt oplyst, før der træffes afgørelse.
  \item[Påstand] Det resultat, en part beder retten om.
  \item[Partshøring] En parts mulighed for at udtale sig om relevante oplysninger, før myndigheden træffer afgørelse.
  \item[Præcedens] En tidligere afgørelse, der har betydning som retligt argument i senere sager.
  \item[Proportionalitet] Kravet om, at et middel skal være egnet, nødvendigt og rimeligt i forhold til formålet.
  \item[Retsfølge] Den juridiske konsekvens, som en regel knytter til bestemte betingelser.
  \item[Retskilde] Et materiale eller princip, der bruges til at fastlægge og begrunde gældende ret.
  \item[Retssubjekt] En person eller organisation, der kan have rettigheder og pligter.
  \item[Skøn] En retlig beslutningsfrihed inden for lovens rammer og formål.
  \item[Subsumption] At sammenholde konkrete fakta med betingelserne i en retlig regel.
  \item[Uagtsomhed] At handle mindre forsigtigt, end situationen kræver, uden nødvendigvis at ønske skaden.
  \item[Ugyldighed] At en retshandel eller afgørelse ikke får den tilsigtede retlige virkning på grund af en retlig fejl.
  \item[Uskyldsformodning] Princippet om, at en person skal behandles som uskyldig, indtil skyld er bevist.
  \item[Gyldighed] Om en handling, aftale, afgørelse eller regel har den tilsigtede retlige virkning.
  \item[Ækvivalens] I EU-retten blandt andet kravet om, at nationale regler for EU-krav ikke må være mindre gunstige end regler for tilsvarende nationale krav.
\end{description}
